\section{Conclusion}
\label{sec:conclusion}

We submitted the four canonical economic cycles---Kitchin, Juglar,
Kuznets and Kondratieff---to a pre-registered three-gate falsifiable
protocol applied across six independent macroeconomic panels spanning
1700--2024. The cycles were armed with four heterogeneous estimators
and tested against a conservative dual null. None of the four
periodicities survives the cascade on more than a small minority of
panels, and Kondratieff is the most fragile: when the per-variable
safeguard of \citet{wen2005} is enforced, zero individual variables
survive on the long-history panels. The few apparent survivals on
income-aggregate composites are diagnosed as compositing artefacts.

The result reproduces, under a falsifiable design, the empirical
consensus of four successive waves of scepticism:
\citet{garvy1943} on Kondratieff, \citet{klotz_neal1973} and
\citet{solomou1987} on Kuznets and Kondratieff, \citet{romer1999} and
\citet{stock_watson2003} on Juglar, and \citet{wen2005} on Kitchin.
The negative result is therefore not a methodological failure---it
is the quantitative confirmation of a long-standing
empirical consensus that has been built piecemeal across the
twentieth century and that, until now, has lacked a unifying
pre-registered protocol.

Two methodological lessons follow. First, claims about canonical
periodicities should be treated as conjectures and tested against
conservative nulls on the strongest data available, not asserted on
the basis of visual inspection of low-pass-filtered composites.
Second, the well-known compositing artefact \citep{wen2005} is the
single largest source of false positives in the literature; a
per-variable safeguard should be a standard requirement.

The substantive question that motivates this paper---\emph{what is
the empirical structure of macroeconomic fluctuations?}---admits an
answer that is statistically robust, theoretically well-founded and
pre-registered, but lies outside the cycle paradigm. It is sketched
in Section~\ref{sec:pointers} and developed in a companion paper.
The Kitchin, Juglar, Kuznets and Kondratieff cycles are not
empirically supported as universal macroeconomic phenomena; the
profession's modelling effort should move beyond them.

\section*{Pre-registration and replication}

The protocol thresholds, cycle bands and surrogate counts used here
were frozen in source code at project inception
(\path{ecowave/cycles/bands.py:CYCLE_BANDS}, immutable
\texttt{MappingProxyType}). The complete pipeline---data ingestion,
Gate~1--3 evaluation, per-variable safeguard---runs end-to-end in a
Docker container in under thirty minutes. The container, the data
manifests with SHA-256 checksums, and the figure-generation scripts
are available on request to the corresponding author and will be
deposited in a Zenodo archive upon journal acceptance.

\section*{Acknowledgements}

This work was developed in the EcoWave research project. The author
thanks the contributors to the J\`ord\`a-Schularick-Taylor
Macrohistory Database, the Bank of England's OBRA team, and the
Bank for International Settlements for the public availability of
the data on which this study rests.

\section*{Declaration of competing interests}

The author declares no competing interests.

\section*{Data availability}

All data are derived from public-domain or academically-licensed
sources listed in Appendix~\ref{app:data} with full provenance and
SHA-256 checksums.
